%% ma: L12-B05-0031
%% lop: 12
%% chuong: 1
%% bai: 5
%% dang: 12.5.22
%% loai: DS
%% muc: [NB, TH, TH, TH]
%% dap_an: "ĐSSS"
%% nguon: "HK1-DeOn_1_1.pdf (M10, Đề ôn tập học kì 1 Toán 12, Đề số 2), Phần II Câu 3"
%% kien_thuc: "Bài 3 (đường tiệm cận), Bài 5 (ứng dụng đạo hàm giải quyết vấn đề thực tiễn)"
%% trang_thai: chinh-thuc
%% ghi_chu: "Đã duyệt 29/10/2026."
\begin{ex}
Một công ty thương mại đưa ra hàm cầu (hàm giá) cho sản phẩm là $p(x)=1600-40x$, trong đó $p(x)$ (đơn vị tính triệu đồng) là giá bán của mỗi sản phẩm mà tại giá bán này có $x$ sản phẩm được bán ra.
\choiceTF{\True Hàm doanh thu của công ty là $R(x)=p(x)\cdot x$}{Hàm doanh thu $R(x)$ có đạo hàm là $R'(x)=80x-1600$}{$R'(x)=0$ có nghiệm là $x=40$}{Đồ thị hàm số $R(x)$ có tiệm cận đứng $x=0$}
\loigiai{a) Đúng.

b) Sai: $R(x)=1600x-40x^2$, $R'(x)=1600-80x$.

c) Sai: $R'(x)=0\Leftrightarrow x=20$.

d) Sai: $R(x)$ là hàm đa thức nên đồ thị không có tiệm cận đứng.}
\end{ex}
